% =====================================================================
% III. JUSTIFICACIÓN
% =====================================================================
\section{Justificaci\'on}

Debido a los obst\'aculos que presenta el an\'alisis de m\'usica escrita a mano se hace necesaria una plataforma de digitalizaci\'on asistida de partituras, justificada en base a la complejidad y variabilidad de la escritura a mano \cite{ricordi2024}. Donde primero, es necesario un an\'alisis sistem\'atico de partituras, para abarcar el mayor n\'umero de variaciones reales y artificiales, buscando la comprensi\'on de distintas escrituras que pueden variar inclusive de naci\'on a naci\'on y a trav\'es de las \'epocas; inclusive, al d\'ia de hoy nuevos s\'imbolos musicales son introducidos en partituras de m\'usica moderna \cite{rebelo2012}.

M\'as all\'a de las variaciones, este proyecto tambi\'en busca colaborar en la detecci\'on de errores. Los errores pueden ser accidentales o inducidos, pero en ambos casos las m\'aquinas tienden a malinterpretar dichos errores. Un ejemplo que ayuda a entender esto es la ``semitonia intelecta'', un problema que conduce a la p\'erdida de elementos musicales en manuscritos antiguos donde los errores no est\'an claramente indicados \cite{rizo2023}. Esto justifica un reconocimiento que vaya m\'as all\'a de la lectura: es necesario, como se ha indicado, el uso de tecnolog\'ias actuales para la detecci\'on de errores, aplicando dentro de este proyecto tecnolog\'ias de machine learning para conseguir un avance significativo en el campo del an\'alisis musical enfocado en la detecci\'on de errores en partituras, un avance que beneficia tanto a los campos dentro del \'ambito musical que buscan la migraci\'on de partituras a formatos digitales como, actualmente, un avance en la ense\~nanza y aprendizaje de escritura musical.
